In the Fokker-Planck equation, the diffusion and convection elements are
bounce-averaged according to the expressions (\ref{diffbav})-(\ref{convbav}%
), which gives, using (\ref{elec}), 
\[
{\Bbb D}_{{\bf p}}^{\text{E}(0)}=0 
\]
and the convection components 
\begin{eqnarray}
F_{p}^{\text{E}(0)} &=&\left\{ q_{e}\xi E_{\Vert }\right\} \\
F_{\xi }^{\text{E}(0)} &=&-\sigma \left\{ \frac{\sigma \xi }{\sqrt{\Psi }\xi
_{0}}q_{e}\sqrt{1-\xi ^{2}}E_{\Vert }\right\}  \nonumber
\end{eqnarray}

Since the poloidal dependence of the electric field on a flux-surface is
given by () 
\begin{equation}
E_{\Vert }\left( \psi ,\theta \right) =\frac{1}{\Psi \left( \psi ,\theta
\right) }\frac{R^{2}\left( \psi ,\theta _{0}\right) }{R^{2}\left( \psi
,\theta \right) }E_{\Vert }^{\left( 0\right) }\left( \psi \right)
\end{equation}
we find 
\begin{eqnarray}
F_{p}^{\text{E}(0)} &=&\frac{\overline{\lambda }_{1,-1,2}}{\lambda }q_{e}\xi
_{0}E_{\Vert }^{\left( 0\right) }\left( \psi \right) \\
F_{\xi }^{\text{E}(0)} &=&-\frac{\overline{\lambda }_{1,-1,2}}{\lambda }q_{e}%
\sqrt{1-\xi _{0}^{2}}E_{\Vert }^{\left( 0\right) }\left( \psi \right) 
\nonumber
\end{eqnarray}
where we defined the bounce averaged coefficient (\ref{lkm}) 
\begin{equation}
\overline{\lambda }_{1,-1,2}=\lambda \sigma \left\{ \sigma \frac{\xi }{\xi
_{0}}\frac{1}{\Psi \left( \psi ,\theta \right) }\frac{R^{2}\left( \psi
,\theta _{0}\right) }{R^{2}\left( \psi ,\theta \right) }\right\}
\end{equation}

\subsubsection{Expression of $\overline{\lambda }_{1,-1,2}$}

This coefficient is sometimes noted $s^{*}$ in the litterature and is
expressed as 
\begin{equation}
\overline{\lambda }_{1,-1,2}=\frac{\sigma }{\widetilde{q}}\left[ \frac{1}{2}%
\sum_{\sigma }\right] _{T}\int_{\theta _{\min }}^{\theta _{\max }}\frac{%
d\theta }{2\pi }\frac{1}{\left| \widehat{\psi }\cdot \widehat{r}\right| }%
\frac{r}{R_{0}}\frac{B}{B_{P}}\sigma \frac{1}{\Psi \left( \psi ,\theta
\right) }\frac{R^{2}\left( \psi ,\theta _{0}\right) }{R^{2}\left( \psi
,\theta \right) }
\end{equation}

Since the integral is odd in $\sigma $, the sum over trapped particles
vanishes, and we have 
\begin{equation}
\overline{\lambda }_{1,-1,2}=\left| 
\begin{array}{ll}
0 & \text{for trapped particles} \\ 
\lambda _{1,-1,2}^{P} & \text{for passing particles}
\end{array}
\right.
\end{equation}
\begin{eqnarray}
\lambda _{1,-1,2}^{P} &=&\frac{1}{\widetilde{q}}\int_{0}^{2\pi }\frac{%
d\theta }{2\pi }\frac{1}{\left| \widehat{\psi }\cdot \widehat{r}\right| }%
\frac{r}{R_{0}}\frac{B}{B_{P}}\frac{1}{\Psi \left( \psi ,\theta \right) }%
\frac{R^{2}\left( \psi ,\theta _{0}\right) }{R^{2}\left( \psi ,\theta
\right) } \\
&=&\frac{1}{\widetilde{q}}\frac{R\left( \psi ,\theta _{0}\right) }{R_{0}}%
\int_{0}^{2\pi }\frac{d\theta }{2\pi }\frac{1}{\left| \widehat{\psi }\cdot 
\widehat{r}\right| }\frac{r}{R\left( \psi ,\theta \right) }\frac{B\left(
\psi ,\theta _{0}\right) }{B_{P}}\frac{R\left( \psi ,\theta _{0}\right) }{%
R\left( \psi ,\theta \right) } \\
&=&\frac{1}{\widetilde{q}}\frac{R\left( \psi ,\theta _{0}\right) }{R_{0}}%
\frac{B\left( \psi ,\theta _{0}\right) }{B_{T}\left( \psi ,\theta
_{0}\right) }\int_{0}^{2\pi }\frac{d\theta }{2\pi }\frac{1}{\left| \widehat{%
\psi }\cdot \widehat{r}\right| }\frac{r}{R\left( \psi ,\theta \right) }\frac{%
B_{T}\left( \psi ,\theta _{0}\right) }{B_{P}}\frac{R\left( \psi ,\theta
_{0}\right) }{R\left( \psi ,\theta \right) } \\
&=&\frac{1}{\widetilde{q}}\frac{R\left( \psi ,\theta _{0}\right) }{R_{0}}%
\frac{B\left( \psi ,\theta _{0}\right) }{B_{T}\left( \psi ,\theta
_{0}\right) }\int_{0}^{2\pi }\frac{d\theta }{2\pi }\frac{1}{\left| \widehat{%
\psi }\cdot \widehat{r}\right| }\frac{r}{R}\frac{B_{T}}{B_{P}} \\
&=&\frac{q\left( \psi \right) }{\widetilde{q}}\frac{R\left( \psi ,\theta
_{0}\right) }{R_{0}}\frac{B\left( \psi ,\theta _{0}\right) }{B_{T}\left(
\psi ,\theta _{0}\right) }
\end{eqnarray}

\subsubsection{Case of circular concentric flux-surfaces}

Note that in the case of circular concentric flux-surfaces, we have 
\begin{eqnarray}
\widetilde{q} &=&\frac{r}{R_{0}}\frac{B\left( \psi ,\theta _{0}\right) }{%
B_{P}\left( \psi ,\theta _{0}\right) } \\
q &=&\frac{r}{R\left( \psi ,\theta _{0}\right) }\frac{B_{T}\left( \psi
,\theta _{0}\right) }{B_{P}\left( \psi ,\theta _{0}\right) }\int_{0}^{2\pi }%
\frac{d\theta }{2\pi }\frac{R\left( \psi ,\theta _{0}\right) }{R}
\end{eqnarray}
so that 
\begin{equation}
\lambda _{1,-1,2}^{P}=\int_{0}^{2\pi }\frac{d\theta }{2\pi }\frac{R\left(
\psi ,\theta _{0}\right) }{R}=\int_{0}^{2\pi }\frac{d\theta }{2\pi }\frac{%
1+\epsilon }{1+\epsilon \cos \theta }
\end{equation}

This integral can then be performed analytically, as done in (\ref{sstar}),
which gives 
\begin{equation}
\lambda _{1,-1,2}^{P}=\sqrt{\frac{1+\epsilon }{1-\epsilon }}
\end{equation}
